% Codex Task 002(B): formal proof of Fourier non-concentration.

\begin{lemma}[Fourier non-concentration]
\label{lem:oracle-comm-fourier-nonconcentration}
Let \(p\ge 7\) be prime, let
\(\mathcal Z_p=\{0\le r<p:b_r\equiv0\pmod p\}\), and put
\[
 F_p(k)=\sum_{r\in\mathcal Z_p}e_p(kr),
 \qquad Z(p)=|\mathcal Z_p|.
\]
For every integer \(1\le K\le(p-1)/2\),
\[
 \sum_{1\le |k|\le K}|F_p(k)|^2
 \le 4\pi^2\left(KZ(p)+\frac{p^2}{K}\right).
\]
In particular the asserted estimate holds with an absolute constant.
\end{lemma}

\begin{proof}
For \(1\le h<p\), write
\[
 A_p(h)=\#\{r:0\le r<r+h<p,\ r,r+h\in\mathcal Z_p\}.
\]
The no-consecutive-zeros lemma gives \(A_p(1)=0\). If \(h\ge2\),
every pair counted by \(A_p(h)\) makes \(r\) a root of the nonzero gap
polynomial \(N_h\bmod p\), of degree \(3(h-1)\). Hence
\begin{equation}
 A_p(h)\le3(h-1).
 \label{eq:oracle-comm-gap-count}
\end{equation}

We first treat \(K\le p/8\). Set
\[
 H=\left\lfloor\frac{p}{4K}\right\rfloor.
\]
Then
\begin{equation}
 \frac{p}{8K}\le H\le\frac{p}{4K},
 \qquad HK/p\le1/4.
 \label{eq:oracle-comm-H-range}
\end{equation}
Partition \(\mathcal Z_p=Z_0\sqcup Z_1\), where
\[
 Z_0\subseteq[0,(p-1)/2],\qquad
 Z_1\subseteq[(p+1)/2,p-1],
\]
and put \(F_i(k)=\sum_{r\in Z_i}e_p(kr)\). Thus
\begin{equation}
 |F_p(k)|^2\le2|F_0(k)|^2+2|F_1(k)|^2.
 \label{eq:oracle-comm-half-cauchy}
\end{equation}
The split into two half intervals is needed to remove the cyclic
wrap-around pairs: within either half, a congruence
\(s-r\equiv h\pmod p\), with \(|h|<H\le p/4\), is the ordinary
integer equality \(s-r=h\).

Use the finite Fej\'er kernel
\[
 \mathcal F_H(k)
 =\frac1H\left|\sum_{u=0}^{H-1}e_p(ku)\right|^2
 =\sum_{|h|<H}\left(1-\frac{|h|}{H}\right)e_p(kh).
\]
For \(1\le |k|\le K\), the sine formula and
\eqref{eq:oracle-comm-H-range} give
\[
 \left|\sum_{u=0}^{H-1}e_p(ku)\right|
 =\frac{|\sin(\pi Hk/p)|}{|\sin(\pi k/p)|}
 \ge\frac{2H}{\pi}.
\]
Indeed, \(\sin(\pi x)\ge2x\) for \(0\le x\le1/4\), while
\(\sin(\pi x)\le\pi x\) for \(x\ge0\). Consequently
\begin{equation}
 \mathbf1_{\{1\le|k|\le K\}}
 \le\frac{\pi^2}{4H}\mathcal F_H(k).
 \label{eq:oracle-comm-fejer-majorant}
\end{equation}

Let \(A_i(h)=\#\{r:r,r+h\in Z_i\}\), with the inequalities interpreted
in the ordinary integer representatives. Finite Fourier orthogonality,
the half-interval observation above, and
\eqref{eq:oracle-comm-fejer-majorant} yield
\begin{align*}
 \sum_{1\le|k|\le K}|F_i(k)|^2
 &\le \frac{\pi^2}{4H}
 \sum_{k\bmod p}\mathcal F_H(k)|F_i(k)|^2\\
 &=\frac{\pi^2p}{4H}
 \left(
 |Z_i|+2\sum_{h=1}^{H-1}
 \left(1-\frac hH\right)A_i(h)
 \right).
\end{align*}
Summing over \(i\), using
\eqref{eq:oracle-comm-half-cauchy}, and observing that
\(A_0(h)+A_1(h)\le A_p(h)\), gives
\begin{align*}
 \sum_{1\le|k|\le K}|F_p(k)|^2
 &\le\frac{\pi^2p}{2H}
 \left(Z(p)+2\sum_{h=1}^{H-1}A_p(h)\right)\\
 &\le\frac{\pi^2p}{2H}
 \left(Z(p)+3H^2\right),
\end{align*}
where the last step follows from
\eqref{eq:oracle-comm-gap-count} and
\[
 \sum_{h=2}^{H-1}3(h-1)
 =\frac32(H-1)(H-2)\le\frac32H^2.
\]
By \eqref{eq:oracle-comm-H-range}, \(p/H\le8K\) and
\(H\le p/(4K)\). Therefore
\[
 \sum_{1\le|k|\le K}|F_p(k)|^2
 \le4\pi^2KZ(p)+\frac{3\pi^2}{4}\frac{p^2}{K}
 \le4\pi^2\left(KZ(p)+\frac{p^2}{K}\right).
\]

It remains to consider \(p/8<K\le(p-1)/2\). Finite Parseval gives
\[
 \sum_{k\bmod p}|F_p(k)|^2=pZ(p).
\]
The residues represented by \(1\le|k|\le K\) are distinct, so their
sum is at most \(pZ(p)\le8KZ(p)\), which is smaller than the asserted
right-hand side. This completes the proof.
\end{proof}

\begin{remark}
If one applies the Fej\'er expansion directly to all of
\(\mathcal Z_p\), its short cyclic correlations also contain ordinary
gaps \(p-h\), for which the bound \(A_p(h)\le3(h-1)\) is unavailable.
The two-half partition above is the required wrap-around correction to
the informal expression
\(2K[Z(p)+\sum_{h\le p/K}A_p(h)(1-hK/p)]\).
\end{remark}
